\section{Finite probability, KL, and contraction calculations}\label{app:probkl}

This appendix records the finite calculations used by the arrow, protocol, and contraction fronts.

\begin{proof}[Proof of Lemma~\ref{lem:deterministic_dpi_kl}]
For $y\in\Omega'$, write $A_y := \{x\in\Omega : g(x)=y\}$ and abbreviate
$p_x := P(x)$, $q_x := Q(x)$, $p_y := \sum_{x\in A_y} p_x$, $q_y := \sum_{x\in A_y} q_x$.
By definition,
\[
D_{\mathrm{KL}}(g_{\#}P \,\|\, g_{\#}Q)
= \sum_{y\in\Omega' : p_y>0} p_y \log\!\Big(\frac{p_y}{q_y}\Big).
\]
For each $y$ with $p_y>0$, the log--sum inequality gives
\[
p_y \log\!\Big(\frac{p_y}{q_y}\Big)
\le \sum_{x\in A_y : p_x>0} p_x \log\!\Big(\frac{p_x}{q_x}\Big),
\]
with the usual convention that the right-hand side is $+\infty$ if some $p_x>0$ but $q_x=0$.
Summing over $y$ and using that the sets $A_y$ partition $\Omega$ yields
\[
D_{\mathrm{KL}}(g_{\#}P \,\|\, g_{\#}Q)
\le \sum_{x\in\Omega : p_x>0} p_x \log\!\Big(\frac{p_x}{q_x}\Big)
= D_{\mathrm{KL}}(P\,\|\,Q).
\]
\end{proof}

\begin{proof}[Proof of Lemma~\ref{lem:rev_pushforward_commute}]
The maps $\mathrm{rev}$ and $f^{(T)}$ commute pointwise:
$f^{(T)}\circ \mathrm{rev} = \mathrm{rev}\circ f^{(T)}$.
Pushforward respects composition, hence
\[
f^{(T)}_{\#}(\mathrm{rev}_{\#}P) = (f^{(T)}\circ \mathrm{rev})_{\#}P
= (\mathrm{rev}\circ f^{(T)})_{\#}P = \mathrm{rev}_{\#}(f^{(T)}_{\#}P).
\]
Substituting $P=P_T^{\mu,K}$ and using the definitions from Section~\ref{sec:setup} gives
\[
P_T^{\mu,K,f,\mathrm{rev}} = \mathrm{rev}_{\#} P_T^{\mu,K,f}
= f^{(T)}_{\#}\big(P_T^{\mu,K,\mathrm{rev}}\big).
\]
\end{proof}

\begin{proof}[Proof of Lemma~\ref{lem:detailed_balance_path}]
Fix $T\ge 0$ and a path $\gamma=(x_0,\dots,x_T)\in X^{T+1}$.
By definition,
\[
P_T^{\pi,K}(\gamma)=\pi(x_0)\prod_{t=0}^{T-1} K(x_t,x_{t+1}),
\qquad
P_T^{\pi,K,\mathrm{rev}}(\gamma)=P_T^{\pi,K}(x_T,\dots,x_0)
=\pi(x_T)\prod_{t=0}^{T-1} K(x_{t+1},x_t).
\]
Using detailed balance repeatedly,
\[
\pi(x_0)K(x_0,x_1)=\pi(x_1)K(x_1,x_0),\;
\pi(x_1)K(x_1,x_2)=\pi(x_2)K(x_2,x_1),\;\dots
\]
Multiplying these identities for $t=0,\dots,T-1$ yields
\[
\pi(x_0)\prod_{t=0}^{T-1}K(x_t,x_{t+1})
=
\pi(x_T)\prod_{t=0}^{T-1}K(x_{t+1},x_t),
\]
which is exactly $P_T^{\pi,K}(\gamma)=P_T^{\pi,K,\mathrm{rev}}(\gamma)$.
Since this holds for every $\gamma$, the two path laws are equal.
\end{proof}

\begin{proof}[Proof of Lemma~\ref{lem:objecthood_dobrushin}]
Fix $A\subseteq Y$ and write $\eta:=\mu-\nu$.
Since $\eta(Y)=0$, if $y_0\in Y$ is arbitrary then
\[
(\mu P)(A)-(\nu P)(A)
=\sum_{y\in Y}\eta(y)P(y,A)
=\sum_{y\in Y}\eta(y)\bigl(P(y,A)-P(y_0,A)\bigr).
\]
Let $Y_+:=\{y\in Y:\eta(y)\ge 0\}$ and $Y_-:=Y\setminus Y_+$.
Because $\sum_y \eta(y)=0$,
\[
\sum_{y\in Y_+}\eta(y)=\sum_{y\in Y_-}(-\eta(y))=d_{\mathrm{TV}}(\mu,\nu).
\]
Hence
\begin{align*}
(\mu P)(A)-(\nu P)(A)
&\le \sum_{y\in Y_+}\eta(y)\sup_{u,v\in Y}\bigl(P(u,A)-P(v,A)\bigr)\\
&\le d_{\mathrm{TV}}(\mu,\nu)\,\sup_{u,v\in Y}\bigl(P(u,A)-P(v,A)\bigr).
\end{align*}
For fixed $u,v$, the variational characterization of total variation gives
\[
P(u,A)-P(v,A)\le d_{\mathrm{TV}}\bigl(P(u,\cdot),P(v,\cdot)\bigr)\le \lambda(P).
\]
Therefore
\[
(\mu P)(A)-(\nu P)(A)\le \lambda(P)\,d_{\mathrm{TV}}(\mu,\nu)
\qquad\text{for every }A\subseteq Y.
\]
Taking the supremum over $A$ yields the stated inequality.
\end{proof}

\begin{proof}[Proof of Lemma~\ref{lem:objecthood_unique_fixed}]
Fix $\mu\in\Delta(Y)$ and set $\mu_t:=\mu P^t$.
By Lemma~\ref{lem:objecthood_dobrushin},
\[
d_{\mathrm{TV}}(\mu_{t+1},\mu_t)
=d_{\mathrm{TV}}(\mu_tP,\mu_{t-1}P)
\le \lambda(P)\,d_{\mathrm{TV}}(\mu_t,\mu_{t-1})
\]
for every $t\ge 1$.
Inducting on $t$ gives
\[
d_{\mathrm{TV}}(\mu_{t+1},\mu_t)\le \lambda(P)^t d_{\mathrm{TV}}(\mu P,\mu).
\]
Hence for integers $m>n$,
\begin{align*}
d_{\mathrm{TV}}(\mu_m,\mu_n)
&\le \sum_{k=n}^{m-1} d_{\mathrm{TV}}(\mu_{k+1},\mu_k)\\
&\le d_{\mathrm{TV}}(\mu P,\mu)\sum_{k=n}^{m-1}\lambda(P)^k.
\end{align*}
Since $0\le \lambda(P)<1$, the geometric tail tends to zero as $n\to\infty$.
Thus $(\mu_t)_{t\ge 0}$ is Cauchy in total variation.
Because $\Delta(Y)$ is a closed subset of the finite-dimensional space $\mathbb{R}^Y$, it is complete for this metric, so there exists $\pi\in\Delta(Y)$ with $\mu_t\to\pi$.
The map $T(\rho)=\rho P$ is continuous, hence
\[
\pi P = \lim_{t\to\infty} \mu_{t+1} = \lim_{t\to\infty}\mu_t = \pi.
\]
So $\pi$ is a fixed distribution.

If $\pi'$ is another fixed distribution, then Lemma~\ref{lem:objecthood_dobrushin} gives
\[
d_{\mathrm{TV}}(\pi,\pi') = d_{\mathrm{TV}}(\pi P,\pi'P)
\le \lambda(P)\,d_{\mathrm{TV}}(\pi,\pi').
\]
Since $\lambda(P)<1$, this forces $d_{\mathrm{TV}}(\pi,\pi')=0$, hence $\pi=\pi'$.
This proves uniqueness.

Finally, applying Lemma~\ref{lem:objecthood_dobrushin} with $\nu=\pi$ and using $\pi P=\pi$ yields
\[
d_{\mathrm{TV}}(\mu P^t,\pi)=d_{\mathrm{TV}}(\mu P^t,\pi P^t)
\le \lambda(P)^t d_{\mathrm{TV}}(\mu,\pi).
\]
\end{proof}

\begin{proof}[Proof of Lemma~\ref{lem:objecthood_approximate_separation}]
By the triangle inequality,
\[
d_{\mathrm{TV}}(\mu,\nu)
\le d_{\mathrm{TV}}(\mu,\mu P)+d_{\mathrm{TV}}(\mu P,\nu P)+d_{\mathrm{TV}}(\nu P,\nu).
\]
Using the hypotheses and Lemma~\ref{lem:objecthood_dobrushin},
\[
d_{\mathrm{TV}}(\mu,\nu)
\le \varepsilon + \lambda(P) d_{\mathrm{TV}}(\mu,\nu) + \varepsilon
\le 2\varepsilon + \lambda d_{\mathrm{TV}}(\mu,\nu).
\]
Rearranging gives
\[
(1-\lambda)d_{\mathrm{TV}}(\mu,\nu)\le 2\varepsilon,
\]
and division by $1-\lambda>0$ proves the claim.
\end{proof}
